\begin{longtable}{L{7cm}rr}

\toprule Показатель & READ COMMITTED & SERIALIZABLE\\\midrule

Целевые запросы, фактически & 1640 & 1070\\

Зафиксированные изменения & 780 & 225\\

Ожидаемые конфликты & 680 & 661\\

Неожиданные отказы & 0 & 4\\

Повторы сериализации & 0 & 517\\

Дубликаты / остатки импорта & 0 / 0 & 0 / 0\\

\bottomrule\end{longtable}

Обе группы использовали один WAR и один восстановленный снимок тестовых таблиц, одинаковый AppState lock, число потоков и предел трёх повторов 40001. При SERIALIZABLE сценарии C05, C10 и C12 остановлены после первого раунда из-за исчерпания повторов и ответа 503; оставшиеся 19 раундов не выполнены и не засчитаны. Отсутствие дубликатов при этом не означает успешность сценария.

Окончательно сохранён READ COMMITTED, заданный явно в бизнес-транзакциях. Глобальная строка-блокировка сериализует окончательные проверки и запись, а ожидаемая версия защищает старую пользовательскую форму. Повышение изоляции при том же глобальном lock не улучшило проверяемые инварианты, но вызвало дополнительные отказы. Короткие транзакции истории также используют READ COMMITTED и не изменяют revision.

Уровень прочитан через SHOW transaction\_isolation в настоящем соединении Database.write и в отдельной транзакции metadata. Ограничение ожидания 200 ms реально дало 55P03 через 217 ms; стандартный Web-порог 10000 ms дал 503 примерно через 10,038 s. Шум общей машины и ранняя остановка не позволяют сравнивать агрегированные времена как строгий benchmark.

Не выполнялась инъекция сетевого разрыва точно в момент ответа commit и нагрузка импортом максимальных 1000 строк. Проверка восстановления RUNNING использует явно созданную изолированную запись и реальный перезапуск единственного процесса; многопроцессное восстановление не заявляется.